% ============================================================
% 论文新章节：ROQ v2 — 旋转 + 双塔 + 自适应 alpha
% ============================================================

\section{From RoPE to OPQ: Rotation-Inspired Quantization}

\subsection{Motivation}

Rotary Position Embedding (RoPE) \cite{rope} 的核心洞察是：
用正交旋转将位置信息编码到相位中，使得在低维子空间中的表示更加紧凑。
这启发了一个自然的问题——能否用类似的旋转变换，
让权重的分布更加"量化友好"？

本章系统性地测试了四种从 RoPE 衍生的思路：
\begin{enumerate}
\item \textbf{激活侧 Hadamard + 权重侧 OPQ}（QuaRot 嫁接）
\item \textbf{可学习旋转矩阵}（训练找最优正交变换）
\item \textbf{Per-channel 自适应 $\alpha$}（用偏度映射）
\item \textbf{双塔量化}（Dual-Tower：大值/小值分离）
\end{enumerate}

\subsection{关键发现 1：Hadamard 直接旋转权重必然失败}

第一轮实验（v1）揭示了一个反直觉的结果：
\textbf{对权重矩阵做 Hadamard 旋转后接 OPQ，SNR 从 15dB 暴跌到 0dB}。

原因分析：
\begin{itemize}
\item Hadamard 变换将能量均匀分布到所有系数上，使得旋转后权重的 std 从 0.02 降至 0
\item OPQ 的 $\alpha < 1$ 变换对小值极其敏感（4th-root 过冲效应）
\item 反旋转 $R @ \hat{W}'$ 将每个元素的量化误差放大 $\sqrt{d}$ 倍
\end{itemize}

这一发现本身具有理论价值：它揭示了 \textbf{OPQ 和 RoPE 的根本差异}——
RoPE 在 head 维度（2D 子空间）独立旋转，不混合跨 head 信息；
而全矩阵 Hadamard 破坏了 OPQ 赖以生存的"小值局部性"。

\subsection{关键发现 2：激活侧旋转 ≠ 权重侧旋转}

修正后的实验（v2）将旋转吸收到激活侧（QuaRot 的正确做法），
但结果仍然是 SNR $\approx$ 0dB。根本原因是：

\textbf{激活侧旋转改变了 $X @ W$ 的有效分布}——
Hadamard 后的激活 $X' = X @ H$ 具有完全不同的统计特性，
使得预先校准的 OPQ scale 失效。

这一发现提示：\textbf{旋转和量化的耦合必须联合校准}，
不能先旋转再独立量化。这是后续可学习旋转方案的设计出发点。

\subsection{关键发现 3：Per-channel 自适应 $\alpha$ 是最强单点改进}

对 512 个输出通道逐一搜索最优 $\alpha$ 后拼接，
4-bit MSE 从 $1.10 \times 10^{-5}$ 降至 $6.46 \times 10^{-6}$，
\textbf{改进 41.2\%}——这是所有方法中最大的单点增益。

更重要的是，我们发现了 $\alpha$ 与通道偏度（skewness）的线性关系：
$$\alpha_i \approx -0.0068 \cdot \text{skew}_i + 0.587$$

用此公式预测 $\alpha$ 后量化，改进 37.1\%，
与逐通道搜索的 41.2\% 仅差 4 个百分点。
这意味着 \textbf{无需搜索，仅通过一次统计计算即可获得近最优的 per-channel $\alpha$}。

\subsection{关键发现 4：双塔量化（Dual-Tower）是全新范式}

受 RoPE "幅值-相位"双路编码启发，我们提出 Dual-Tower Quantization：
将权重按幅值拆分为"大值塔"和"小值塔"，分别用不同的 $\alpha$ 量化。

\begin{equation}
\hat{W} = \underbrace{Q_{\alpha=0.5}(W \odot \mathbb{I}[|W| \geq \tau])}_{\text{Tower 1: Large values (Sqrt)}}
       + \underbrace{Q_{\alpha=0.25}(W \odot \mathbb{I}[|W| < \tau])}_{\text{Tower 2: Small values (4th Root)}}
\end{equation}

实验表明：
\begin{itemize}
\item 最优阈值 $\tau^* \approx 0.015$（约 33\% 大值 + 67\% 小值）
\item 4-bit MSE 改进 22.7\%，3-bit 改进 5.1\%
\item 存储开销：每权重多 1 bit 的 tower 标记位，但 $\alpha$ 差异带来的精度收益远超成本
\end{itemize}

\subsection{最终排名}

\begin{table}[h]
\centering
\caption{4-bit OPQ 增强方法排名（attention\_q 层, dim=512）}
\label{tab:roq_v2_ranking}
\begin{tabular}{c|l|c|c|c}
\hline
Rank & Method & MSE & SNR (dB) & Improvement \\
\hline
1 & Per-channel $\alpha$ search & $6.46 \times 10^{-6}$ & 18.17 & \textbf{+41.2\%} \\
2 & Per-channel $\alpha$ (skew pred) & $6.92 \times 10^{-6}$ & 17.87 & +37.1\% \\
3 & Dual-Tower (optimal $\tau$) & $8.75 \times 10^{-6}$ & 16.85 & +20.4\% \\
4 & Learned Rotation (64d) & $1.06 \times 10^{-5}$ & 16.19 & +3.9\% \\
5 & OPQ $\alpha=0.45$ (baseline) & $1.10 \times 10^{-5}$ & 15.86 & 0.0\% \\
\hline
\end{tabular}
\end{table}

\subsection{理论贡献}

本文从 RoPE 的旋转编码思想出发，最终提炼出三个可独立贡献的理论点：

\begin{enumerate}
\item \textbf{OPQ 的"旋转不兼容性"定理}：
  全矩阵正交旋转与幂次量化不兼容，因为旋转破坏小值局部性。
  这解释了为何 QuaRot 只对激活做旋转，而不对权重做旋转。

\item \textbf{Skewness-$\alpha$ 映射公式}：
  给出了 $\alpha_i = -0.0068 \cdot \text{skew}_i + 0.587$ 的闭式近似，
  使得 per-channel 自适应 OPQ 无需搜索，O(d) 复杂度即可完成。

\item \textbf{Dual-Tower 量化框架}：
  首次提出"按幅值分离 + 异质 $\alpha$"的量化范式，
  为大值和小值分别定制最优变换，开辟了"混合幂次量化"的新方向。
\end{enumerate}

\subsection{局限与未来工作}

\begin{itemize}
\item \textbf{可学习旋转的收敛问题}：当前实验在 64 维下仅改进 2.7\%，
  因为梯度信号被量化阶梯函数淹没。未来可探索 straight-through estimator (STE) 或 Gumbel-softmax 松弛。
\item \textbf{Dual-Tower 的标记位开销}：当前需要 1 extra bit 标记 tower 归属，
  未来可探索"隐式分离"（如利用量化整数的最高位作为隐式标记）。
\item \textbf{跨层联合优化}：per-channel $\alpha$ 和 Dual-Tower $\tau$ 目前是逐层独立优化的，
  未来可探索跨层的联合搜索（类似逐层混合精度的思路）。
\end{itemize}
